\documentclass[11pt]{article}
\usepackage[margin=1in]{geometry}
\usepackage{amsmath,amssymb,amsthm,mathtools,bm}
\usepackage{booktabs,longtable,array}
\usepackage[dvipsnames]{xcolor}
\usepackage[colorlinks=true,linkcolor=MidnightBlue,citecolor=BrickRed,urlcolor=MidnightBlue]{hyperref}
\usepackage{microtype}

\hypersetup{
  pdftitle={One-Step Generated Effective-Action Transfer},
  pdfauthor={},
  pdfsubject={A finite-dimensional retained/fiber RG theorem},
  pdfkeywords={Yang--Mills, renormalization group, effective action, Brascamp--Lieb, fiber marginalization}
}

\newtheorem{theorem}{Theorem}[section]
\newtheorem{proposition}[theorem]{Proposition}
\newtheorem{lemma}[theorem]{Lemma}
\newtheorem{corollary}[theorem]{Corollary}
\theoremstyle{definition}
\newtheorem{definition}[theorem]{Definition}
\newtheorem{remark}[theorem]{Remark}

\newcommand{\R}{\mathbb R}
\newcommand{\E}{\mathbb E}
\newcommand{\Cov}{\operatorname{Cov}}
\newcommand{\Hess}{\operatorname{Hess}}
\newcommand{\Id}{\mathrm I}
\newcommand{\good}{\mathrm g}
\newcommand{\bad}{\mathrm b}
\newcommand{\eff}{\mathrm{eff}}
\newcommand{\statusproved}{\textcolor{ForestGreen}{\textsc{proved}}}
\newcommand{\statusnumeric}{\textcolor{RoyalBlue}{\textsc{computed}}}
\newcommand{\statusopen}{\textcolor{BrickRed}{\textsc{open}}}

\title{One-Step Generated Effective-Action Transfer\\
\large Explicit retained/fiber recursion, centered transport, and defect sensitivity}
\author{Research dossier---restricted theorem and frontier audit}
\date{20 July 2026}

\begin{document}
\maketitle

\begin{abstract}
An explicit four-dimensional cubical block supplies a retained space of
dimension $51$, a conditional $\mathfrak{su}(2)$ fiber of dimension $357$, and
a certified tangent fiber lower bound $\lambda_F=2$.  The next question is
whether a local nonlinear certificate can survive the operation that actually
generates an effective action: integration of the fiber.  This paper proves a
finite-dimensional one-step theorem.  A five-component scaled derivative norm
controls fiber coercivity, conditional-minimizer motion, centered conditional
transport, and the Hessian of the exactly marginalized action.  The generated
good-action budget is
\[
 \eta_{\good}=\varepsilon_Y+
 \frac{\varepsilon_X^2}{\lambda_F-\varepsilon_F}.
\]
A gauge-invariant bad/good ratio with scaled $C^2$ budgets contributes at most
$\zeta_2+\zeta_1^2$.  State normalization and conditional observable transport
are exact, and residual orthogonal covariance is preserved.  These are
\statusproved\ implications from declared finite-chart hypotheses.  Explicit
Yang--Mills initialization of those hypotheses, dependent large-field control,
locality across blocks, repeated ultraviolet iteration, and infrared mass
generation remain \statusopen.
\end{abstract}

\tableofcontents

\section{Target and epistemic boundary}

The preceding quotient/fiber construction isolates two kinds of information.
The retained variable $y$ carries the coarse curvature data that must survive
blocking; the fiber variable $z$ carries the fluctuation information integrated
out at one step.  A local Hessian estimate on $z$ is useful only if it can be
converted into a controlled statement about the generated action on $y$.

The target here is therefore neither a formal beta function nor an infrared
claim.  It is the exact finite-dimensional implication
\[
 \text{input derivative certificate}
 \longrightarrow
 \text{fiber marginalization}
 \longrightarrow
 \text{output derivative certificate}.
\]
Balaban's constructive programme generates localized effective actions in a
small-field approximation through one-step transformations
\cite{Balaban1987}; the theorem below abstracts one analytic mechanism needed
inside such a step.  The heat-kernel action motivating the initialization was
introduced by Menotti and Onofri \cite{MenottiOnofri}.

The theorem is a restricted-model result at claim level C2.  It does not show
that the four-dimensional nonlinear $SU(2)$ measure satisfies its input norm,
that the exceptional set is summable, or that the recursion closes through a
number of scales diverging with the cutoff.  The Clay target remains the
construction of nontrivial pure Yang--Mills theory on $\R^4$ with a positive
Hamiltonian gap \cite{JaffeWitten}.

\section{Typed one-step package}

Let $Y=\R^m$ be the retained space and $F=\R^n$ the conditional fiber.  For the
block imported from the previous stage, $(m,n)=(51,357)$.  Let $D\subset F$ be
a fixed convex domain and $t>0$ the declared heat-time normalization.  Write
\begin{equation}\label{eq:action}
 S(y,z)=\frac{1}{2t}\bigl(\langle y,L_Yy\rangle+
 \langle z,L_Fz\rangle\bigr)+R(y,z),
\end{equation}
where
\begin{equation}\label{eq:tangent-bounds}
 L_F\succeq\lambda_F\Id,
 \qquad L_Y\succeq\lambda_Y\Id.
\end{equation}
The coordinates are Gaussian conditional-mean coordinates: the tangent
quadratic form has no retained--fiber mixed term.  This is a choice of local
splitting, not a global gauge fixing theorem.

Define the good partition function, effective action, and conditional law by
\begin{equation}\label{eq:marginal}
 Z_{\good}(y)=\int_D e^{-S(y,z)}\,dz,
 \quad
 \widehat S_{\good}(y)=-\log Z_{\good}(y),
 \quad
 \nu_y(dz)=Z_{\good}(y)^{-1}e^{-S(y,z)}\,dz.
\end{equation}
All differentiations below are justified either on a bounded $D$ or by the
uniform fiber convexity and the standard domination hypotheses stated with the
theorem.

\begin{definition}[Scaled local action norm]\label{def:norm}
The remainder $R$ has budget
\[
 \mathfrak N_t(R)=
 (\varepsilon_F,\varepsilon_X,\varepsilon_Y,\tau_F,\tau_X)
\]
on a declared region if
\begin{equation}\label{eq:second}
 t\|R_{zz}\|\le\varepsilon_F,\qquad
 t\|R_{zy}\|\le\varepsilon_X,\qquad
 t\|R_{yy}\|\le\varepsilon_Y.
\end{equation}
\begin{equation}\label{eq:third}
 t\|D_zR_{zz}\|\le\tau_F,\qquad
 t\|D_zR_{zy}\|\le\tau_X.
\end{equation}
All norms are Euclidean operator or multilinear operator norms.  The five
entries are dimensionless and are part of the RG state; they are not hidden in
big-$O$ notation.
\end{definition}

\section{Fiber margin and centered transport}

Set
\begin{equation}\label{eq:mu}
 \mu:=\lambda_F-\varepsilon_F.
\end{equation}

\begin{proposition}[Fiber coercivity and minimizer response, \statusproved]
\label{prop:minimizer}
Assume $\mu>0$.  Then
\[
 S_{zz}(y,z)\succeq\frac{\mu}{t}\Id.
\]
Whenever the conditional minimizer $z_\star(y)$ is interior, it is unique and
\begin{equation}\label{eq:response}
 \|Dz_\star(y)\|\le h,
 \qquad h:=\frac{\varepsilon_X}{\mu}.
\end{equation}
\end{proposition}
\begin{proof}
Equations \eqref{eq:tangent-bounds} and \eqref{eq:second} give the Hessian
bound.  Differentiate $S_z(y,z_\star(y))=0$ to obtain
\[
 Dz_\star=-S_{zz}^{-1}S_{zy}.
\]
The tangent mixed Hessian vanishes, so $S_{zy}=R_{zy}$.  Hence
$\|S_{zz}^{-1}\|\le t/\mu$ and $\|S_{zy}\|\le\varepsilon_X/t$, proving
\eqref{eq:response}.
\end{proof}

For $w=z-z_\star(y)$, define the centered potential
\[
 V_y(w)=S(y,z_\star(y)+w)-S(y,z_\star(y)).
\]

\begin{proposition}[Explicit centered cross budget, \statusproved]
\label{prop:cross}
Suppose \eqref{eq:third} holds on every segment
$z_\star(y)+sw$, $0\le s\le1$, with $\|w\|\le r$.  Then
\begin{equation}\label{eq:chi}
 t\|D_y\nabla_wV_y(w)\|
 \le \chi(r):=r\left(\tau_X+\tau_F
 \frac{\varepsilon_X}{\mu}\right).
\end{equation}
\end{proposition}
\begin{proof}
Differentiating $S_z(y,z_\star(y))=0$ also shows that
$D_y\nabla_wV_y(0)=0$.  At a general $w$,
\[
 D_y\nabla_wV_y(w)=S_{zy}(y,z_\star+w)
 +S_{zz}(y,z_\star+w)Dz_\star.
\]
Subtract its value at $w=0$.  The constant tangent fiber Hessian cancels.
The mean-value theorem, \eqref{eq:third}, and \eqref{eq:response} give
\eqref{eq:chi}.
\end{proof}

Let $B_r\subset F$ be a common centered convex tube, and normalize
$e^{-V_y(w)}\,dw$ on $B_r$.

\begin{corollary}[Centered-fiber transport, \statusproved]
\label{cor:transport}
For retained points connected by a segment inside the declared chart,
\begin{equation}\label{eq:K}
 W_2(\widetilde\nu_y,\widetilde\nu_{y'})
 \le K_F(r)\|y-y'\|,
 \qquad
 K_F(r):=\frac{\chi(r)}{\mu}.
\end{equation}
Thus $K_F(r)<1$ is an explicit local centered-fiber contraction certificate.
\end{corollary}
\begin{proof}
The centered potentials are $\mu/t$ strongly convex.  The
Brascamp--Lieb/logarithmic-Sobolev transport argument bounds response by the
ratio of the cross-gradient constant $\chi(r)/t$ to the convexity constant
$\mu/t$; see also \cite{OttoVillani}.  Equation \eqref{eq:K} follows.
\end{proof}

\section{Exact marginalization and the generated Hessian}

\begin{theorem}[One-step generated-action bound, \statusproved]
\label{thm:generated}
Assume $R$ is $C^2$, $\mu>0$, differentiation under
\eqref{eq:marginal} is valid, and the Brascamp--Lieb covariance inequality
applies to $\nu_y$.  Then
\begin{equation}\label{eq:loewner}
 -\left(\varepsilon_Y+\frac{\varepsilon_X^2}{\mu}\right)\Id
 \preceq
 t\nabla_y^2\widehat S_{\good}(y)-L_Y
 \preceq \varepsilon_Y\Id.
\end{equation}
In particular,
\begin{equation}\label{eq:eta-good}
 \left\|t\nabla_y^2\widehat S_{\good}(y)-L_Y\right\|
 \le \eta_{\good},
 \qquad
 \eta_{\good}:=\varepsilon_Y+
 \frac{\varepsilon_X^2}{\mu}.
\end{equation}
\end{theorem}
\begin{proof}
Twice differentiating the logarithm of the partition function gives the exact
identity
\begin{equation}\label{eq:hessian-identity}
 \nabla_y^2\widehat S_{\good}
 =\E_{\nu_y}[S_{yy}]-\Cov_{\nu_y}(S_y,S_y).
\end{equation}
The deterministic part $t^{-1}L_Yy$ has zero covariance.  For every unit
vector $u\in Y$, the Brascamp--Lieb covariance inequality
\cite{BrascampLieb} gives
\begin{align*}
 \operatorname{Var}_{\nu_y}(u\cdot R_y)
 &\le \E_{\nu_y}\!\left[
  \langle R_{zy}u,S_{zz}^{-1}R_{zy}u\rangle\right]\\
 &\le \frac{\varepsilon_X^2}{t\mu}.
\end{align*}
Moreover, $-\varepsilon_Y\Id/t\preceq R_{yy}\preceq
\varepsilon_Y\Id/t$.  Insert these estimates in
\eqref{eq:hessian-identity}.  Covariance is positive semidefinite, yielding
both sides of \eqref{eq:loewner} and hence \eqref{eq:eta-good}.
\end{proof}

The term $\varepsilon_X^2/\mu$ is the nonlinear analogue of a Schur-complement
correction.  In a Gaussian model, the generated Hessian is exactly
\[
 A_{\eff}=A-BC^{-1}B^*.
\]
The accompanying implementation checks this identity and its covariance under
orthogonal retained and fiber recodings.

\section{The defect channel}

Let the full one-block partition function split into good and bad parts:
\[
 Z_{\mathrm{full}}(y)=Z_{\good}(y)+Z_{\bad}(y)
 =Z_{\good}(y)(1+\delta(y)),
 \qquad \delta(y):=\frac{Z_{\bad}(y)}{Z_{\good}(y)}\ge0.
\]
Consequently,
\begin{equation}\label{eq:defect-action}
 \widehat S_{\mathrm{full}}(y)
 =\widehat S_{\good}(y)-\log(1+\delta(y)).
\end{equation}

\begin{definition}[Scaled defect norm]
On the retained chart, suppose
\begin{equation}\label{eq:defect-norm}
 0\le\delta\le\zeta_0,
 \qquad \sqrt t\|D\delta\|\le\zeta_1,
 \qquad t\|D^2\delta\|\le\zeta_2.
\end{equation}
The activity $\zeta_0$ controls mass comparison, whereas
$\zeta_1,\zeta_2$ control the generated derivative budget.
\end{definition}

\begin{proposition}[Defect sensitivity, \statusproved]
\label{prop:defect}
Under \eqref{eq:defect-norm},
\begin{equation}\label{eq:defect-hessian}
 t\left\|D^2\log(1+\delta)\right\|
 \le \eta_{\bad}:=\zeta_2+\zeta_1^2.
\end{equation}
Therefore
\begin{equation}\label{eq:eta-full}
 \left\|t\nabla_y^2\widehat S_{\mathrm{full}}-L_Y\right\|
 \le
 \eta_{\mathrm{full}}:=
 \varepsilon_Y+
 \frac{\varepsilon_X^2}{\lambda_F-\varepsilon_F}
 +\zeta_2+\zeta_1^2.
\end{equation}
If $L_Y\succeq\lambda_Y\Id$ and
$\eta_{\mathrm{full}}<\lambda_Y$, the generated base action is at least
$(\lambda_Y-\eta_{\mathrm{full}})/t$ strongly convex on this finite chart.
\end{proposition}
\begin{proof}
Differentiate $\log(1+\delta)$ twice:
\[
 D^2\log(1+\delta)=\frac{D^2\delta}{1+\delta}
 -\frac{D\delta\otimes D\delta}{(1+\delta)^2}.
\]
Since $\delta\ge0$, \eqref{eq:defect-hessian} follows from the triangle
inequality.  Combine it with Theorem~\ref{thm:generated} and
\eqref{eq:defect-action}.
\end{proof}

The proposition identifies exactly what a large-field theorem must supply.
It does not prove the bounds \eqref{eq:defect-norm}.  In particular,
$\zeta_0$ must eventually arise from a dependent polymer estimate rather than
from treating blocks as independent.

\section{State, observables, and covariance}

For an integrable observable $O(y,z)$, define its blocked representative
\begin{equation}\label{eq:observable}
 (\mathcal TO)(y)=\int_D O(y,z)\,\nu_y(dz).
\end{equation}
Fubini's theorem gives the exact state-transport identity
\begin{equation}\label{eq:state}
 \int O(y,z)e^{-S(y,z)}\,dy\,dz
 =\int (\mathcal TO)(y)e^{-\widehat S_{\good}(y)}\,dy.
\end{equation}
Thus marginalization preserves positivity and normalization and transports
observables by conditional expectation.  It is lossy as a map of individual
configurations but exact as a state/observable duality.

Let a compact residual group $H$ act orthogonally by
$(y,z)\mapsto(\rho_Y(h)y,\rho_F(h)z)$.  If $D$ is invariant and
$S(\rho_Y(h)y,\rho_F(h)z)=S(y,z)$, then a change of variables proves
\[
 \widehat S_{\good}(\rho_Y(h)y)=\widehat S_{\good}(y).
\]
Equivariant observables remain equivariant under $\mathcal T$.  This is the
covariance of the local residual representation.  A global gauge-orbit atlas,
including exceptional stabilizers, is a separate open problem.

\section{RG certificate and parameter registry}

\begin{longtable}{@{}p{0.22\textwidth}p{0.30\textwidth}p{0.38\textwidth}@{}}
\toprule
Certificate field & Stage-06 result & Qualification\\
\midrule
\endhead
Configuration spaces & finite $Y\oplus F$ & local quotient/fiber chart\\
Blocking kernel & exact conditional integration & fixed declared domain\\
Gauge covariance & preserved for invariant orthogonal residual action & global orbit atlas open\\
State transport & exact, Eq.~\eqref{eq:state} & good channel; defect added by Eq.~\eqref{eq:defect-action}\\
Observable transport & exact conditional expectation & support/locality growth not bounded\\
Measure positivity & preserved & weaker than OS reflection positivity\\
Physical scale & $t$ explicit; block factor must be declared externally & no energy-unit mass conclusion\\
Locality & absent from the finite norm & polymer interaction norm open\\
Volume dependence & one finite block & uniform-volume theorem open\\
Running coupling & not extracted & initialization and scheme open\\
Reconstruction & conditional law $\nu_y$ retained & individual fiber configuration forgotten\\
Error recursion & Eqs.~\eqref{eq:K}, \eqref{eq:eta-good}, \eqref{eq:eta-full} & one step only\\
\bottomrule
\end{longtable}

\begin{center}
\begin{tabular}{@{}lll@{}}
\toprule
Parameter & Meaning & Required condition\\
\midrule
$\lambda_F$ & tangent fiber lower bound & imported value $2$\\
$\varepsilon_F$ & nonlinear fiber-Hessian budget & $\varepsilon_F<\lambda_F$\\
$\varepsilon_X$ & retained--fiber mixed budget & finite\\
$\varepsilon_Y$ & direct retained-Hessian budget & finite\\
$\tau_F,\tau_X$ & fiber variation of second derivatives & finite on declared tube\\
$r$ & centered good-tube radius & $K_F(r)<1$ for contraction\\
$\zeta_0,\zeta_1,\zeta_2$ & defect mass and derivative budgets & must be proved from full measure\\
$\lambda_Y$ & finite-chart base reference curvature & not a physical gap\\
\bottomrule
\end{tabular}
\end{center}

\section{Executable calibration and hostile cases}

The implementation evaluates the formulas without asymptotic notation.  With
$\lambda_F=2$, $\lambda_Y=1$, and the small example
\[
 (\varepsilon_F,\varepsilon_X,\varepsilon_Y,\tau_F,\tau_X)
 =(0.25,0.20,0.10,0.30,0.20),
\]
$r=0.5$, and $(\zeta_0,\zeta_1,\zeta_2)=(0.01,0.02,0.03)$, it returns
\[
 \mu=1.75,
 \quad K_F\approx0.06694,
 \quad \eta_{\good}\approx0.12286,
 \quad \eta_{\mathrm{full}}\approx0.15326.
\]
A deterministic $3+4$ dimensional Gaussian witness has actual generated
remainder norm $0.0937808$, below the certified bound $0.107226$.

The test suite also enforces the boundary of the theorem:
\begin{itemize}
\item $\varepsilon_F\ge\lambda_F$ fails closed;
\item a large tube may retain fiber coercivity while losing $K_F<1$;
\item a nonpositive fiber Hessian and a nonorthogonal claimed frame are rejected;
\item $\lambda_Y=0$ is never converted into a positive mass claim;
\item the Gaussian Schur complement is invariant under orthogonal retained and
fiber recodings.
\end{itemize}

These tests mechanize finite algebra only.  They do not sample gauge fields or
certify the analytic initialization.

\section{Frontier audit and restrategy}

\begin{longtable}{@{}p{0.08\textwidth}p{0.66\textwidth}p{0.14\textwidth}@{}}
\toprule
ID & Required implication & Status\\
\midrule
\endhead
R1 & One-step derivative, transport, and defect recursion from declared norms & \statusproved\\
F1 & Explicit cutoff-aware initialization of all five derivative budgets for the exact nonlinear $SU(2)$ heat-kernel block and generated interaction & \statusopen\\
F2 & Gauge-invariant dependent polymer estimate proving the defect budgets, including topology and exceptional orbits & \statusopen\\
F3 & Uniform spatial locality, observable-support control, and positivity preservation or recovery & \statusopen\\
F4 & Closure of the recursion for $k(a)\to\infty$ steps with summable errors & \statusopen\\
F5 & Nonperturbative crossover of the retained base into a massive basin & \statusopen\\
F6 & Nontrivial regulator-independent OS continuum theory and Hamiltonian spectral transfer & \statusopen\\
\bottomrule
\end{longtable}

The strongest advance is representational: the local fiber certificate is now
connected to the exact operation that forgets the fiber and generates the next
description.  The theorem explains which information survives and the exact
price of forgetting it.  It also localizes the first unproved implication.
The obstruction is no longer the algebra of marginalization; it is proving
that the nonlinear gauge theory enters the declared norm, including its bad
channel, with constants stable in cutoff and volume.

The next target should therefore be an initialization theorem, not another
abstract recursion.  It should compute explicit $SU(2)$ heat-kernel/BCH chart
constants for the five norm components and pair them with a genuinely
gauge-invariant defect activity.  If the resulting values violate
$\varepsilon_F<2$ or $K_F<1$, the failure will identify whether the retained
variables, chart, or good/bad decomposition must change.

\begin{thebibliography}{9}
\bibitem{Balaban1987}
T.~Balaban,
\emph{Renormalization group approach to lattice gauge field theories: I.
Generation of effective actions in a small field approximation and a coupling
constant renormalization in four dimensions},
Commun. Math. Phys. \textbf{109} (1987), 249--301,
\url{https://doi.org/10.1007/BF01215223}.

\bibitem{BrascampLieb}
H.~J.~Brascamp and E.~H.~Lieb,
\emph{On extensions of the Brunn--Minkowski and Pr\'ekopa--Leindler theorems,
including inequalities for log concave functions, and with an application to
the diffusion equation},
J. Funct. Anal. \textbf{22} (1976), 366--389,
\url{https://doi.org/10.1016/0022-1236(76)90004-5}.

\bibitem{MenottiOnofri}
P.~Menotti and E.~Onofri,
\emph{The action of $SU(N)$ lattice gauge theory in terms of the heat kernel on
the group manifold},
Nucl. Phys. B \textbf{190} (1981), 288--300,
\url{https://doi.org/10.1016/0550-3213(81)90560-5}.

\bibitem{OttoVillani}
F.~Otto and C.~Villani,
\emph{Generalization of an inequality by Talagrand and links with the
logarithmic Sobolev inequality},
J. Funct. Anal. \textbf{173} (2000), 361--400,
\url{https://doi.org/10.1006/jfan.1999.3557}.

\bibitem{JaffeWitten}
A.~Jaffe and E.~Witten,
\emph{Quantum Yang--Mills Theory}, Clay Mathematics Institute Millennium
Problem description,
\url{https://www.claymath.org/wp-content/uploads/2022/06/yangmills.pdf}.
\end{thebibliography}

\end{document}
